
\subsubsection{Mechanism Validation Interpretation}

The 100\% gate pass rate provides evidence that feedback modality curriculum is implementable. Phase 1 execution weight (0.$800\rightarrow 0.714)$ establishes correctness foundation, Phase 2 AI weight peak (0.545 at 50\%) enables scalable quality refinement, and Phase 3 human weight increase (0.$400\rightarrow 0.636)$ prevents execution-only collapse. Each transition occurs smoothly through our programmed Gaussian and linear schedules.

The sequential validation through prerequisite chain (h-e1 $\rightarrow$ h-m1 $\rightarrow$ h-m2 $\rightarrow$ h-m3) demonstrates later phases build on earlier foundations. Phase 2 quality improvement would not be meaningful without Phase 1 correctness baseline. Phase 3 conflict case analysis requires Phase 2 to have established which samples are edge cases.

The dual improvement in Phase 2---where both quality and correctness increase simultaneously---challenges the assumption that feedback signals capture fully orthogonal quality dimensions. This suggests AI reward models learn representations that partially overlap with both signals. Importantly, this does not invalidate multi-modal integration: even partially orthogonal signals add unique information.

\subsubsection{Limitations}

Three limitations constrain claims we can make, disclosed transparently.

\textbf{Limitation 1: Performance Untested.} All experiments used pretrained CodeGen-350M without reinforcement learning training. We simulated training trajectories through checkpoint evaluation rather than actual policy gradient updates and reward optimization. Consequently, we validate mechanism feasibility---the framework implements correctly and predicted patterns emerge---not quantitative performance gains. Whether mechanism advantages translate to performance improvements over baselines requires full-scale RL training with reward optimization. Future work requires actual RL training with PPO, comparison against baselines trained to convergence, multiple random seeds, and evaluation on independent test sets.

\textbf{Limitation 2: Heuristic Human Feedback.} Human preference scores use code quality indicators (length, documentation, structural patterns) rather than actual human annotator ratings. This affects external validity but not mechanism validation---the weight scheduling logic is orthogonal to the quality of the human signal. Real annotations would strengthen external validity but would not change the mechanism validation conclusion. Future work should collect annotated samples with inter-annotator agreement analysis and retrain AI reward models on real preference pairs.

\textbf{Limitation 3: No Static Comparison.} We compare tri-modal dynamic scheduling against single-feedback baselines conceptually but not against tri-modal static (fixed optimal weights throughout training). This means we cannot claim dynamic scheduling outperforms static integration. Our contribution is demonstrating dynamic scheduling is viable---it can be implemented, produces predicted patterns, and achieves phase-specific objectives. Whether dynamic is optimal compared to static remains open. Future work should implement grid search over static weight configurations and compare against dynamic schedules.

\subsubsection{Broader Implications}

Our work establishes feedback modality scheduling as a research direction for multi-objective RL. Existing curriculum learning schedules task difficulty or model capacity. We schedule which feedback signal to emphasize, opening design space: when should correctness signal dominate? when do learned reward models become critical? when is human oversight most valuable? These questions generalize beyond code generation to any domain with heterogeneous feedback sources.

The mechanism validation approach---testing predicted patterns through gate criteria rather than maximizing benchmark performance---reflects methodological choice. Benchmark-driven research optimizes for state-of-the-art numbers, often conflating mechanism novelty with hyperparameter tuning. Mechanism-first validation separates these concerns: we establish the proposed mechanism works as designed, then leave performance optimization to follow-up work.

\subsubsection{Future Directions}

Immediate extensions include: (1) full RL training validation with performance comparison against baselines, (2) static vs. dynamic ablation to test optimality claims, (3) real human annotation collection with inter-rater reliability analysis, (4) fourth feedback modality integration (static analysis tools), and (5) alternative schedule parameterizations (learned schedules, adaptive schedules).

Medium-term extensions could explore: (1) multi-file code generation on repository-level tasks (does feedback curriculum generalize beyond competitive programming?), (2) cross-domain application (image generation, dialogue systems), and (3) theoretical analysis of convergence properties.

Long-term vision includes meta-learning optimal schedules automatically across tasks, learned weight schedules that adapt to training dynamics, and unified theory of curriculum learning encompassing task difficulty, model capacity, and feedback modality dimensions.
